\documentclass[a4paper,11pt]{book}
\usepackage{listings}
\usepackage[utf8]{inputenc}
\usepackage[spanish]{babel}

\decimalpoint
\usepackage{dcolumn}
\newcolumntype{.}{D{.}{\esperiod}{-1}}
\makeatletter
\addto\shorthandsspanish{\let\esperiod\es@period@code}
\makeatother

\RequirePackage{verbatim}
\usepackage{fancyhdr}
\usepackage{graphicx}
\usepackage{afterpage}

\usepackage{longtable}

\usepackage[pdfborder={000}]{hyperref} %referencia

% ********************************************************************
% Re-usable information
% ********************************************************************
\newcommand{\myDocument}{Memoria de prácticas}
\newcommand{\myTitle}{Prácticas de empresa}
\newcommand{\myCompany}{Bi 4 Group}
\newcommand{\myDegree}{Máster en Ingeniería Informática}
\newcommand{\myName}{Aythami Estévez Olivas}
\newcommand{\myProf}{Pedro Villar Castro (tutor académico)}
\newcommand{\myOtherProf}{Nombre Apllido1 Apellido2 (tutor en la empresa)}
\newcommand{\myFaculty}{Escuela Técnica Superior de Ingenierías Informática y de Telecomunicación}
\newcommand{\myFacultyShort}{E.T.S.I.I.T.}
\newcommand{\myUni}{\protect{Universidad de Granada}}
\newcommand{\myLocation}{Granada}
\newcommand{\myTime}{\today}


\hypersetup{
pdfauthor = {\myName (aythae@correo.ugr.es)},
pdftitle = {\myTitle},
pdfsubject = {},
pdfkeywords = {prácticas_empresa, memoria, bi4group},
pdfcreator = {LaTeX con el paquete MiKTeX},
pdfproducer = {pdflatex}
}

\usepackage{url}
\usepackage{colortbl,longtable}
\usepackage[stable]{footmisc}

\pagestyle{fancy}
\fancyhf{}
\fancyhead[LO]{\leftmark}
\fancyhead[RE]{\rightmark}
\fancyhead[RO,LE]{\textbf{\thepage}}
\renewcommand{\chaptermark}[1]{\markboth{\textbf{#1}}{}}
\renewcommand{\sectionmark}[1]{\markright{\textbf{\thesection. #1}}}

\setlength{\headheight}{1.5\headheight}

\newcommand{\HRule}{\rule{\linewidth}{0.5mm}}

\definecolor{gray97}{gray}{.97}
\definecolor{gray75}{gray}{.75}
\definecolor{gray45}{gray}{.45}
\definecolor{gray30}{gray}{.94}

\lstset{ frame=Ltb,
     framerule=0.5pt,
     aboveskip=0.5cm,
     framextopmargin=3pt,
     framexbottommargin=3pt,
     framexleftmargin=0.1cm,
     framesep=0pt,
     rulesep=.4pt,
     backgroundcolor=\color{gray97},
     rulesepcolor=\color{black},
     %
     stringstyle=\ttfamily,
     showstringspaces = false,
     basicstyle=\scriptsize\ttfamily,
     commentstyle=\color{gray45},
     keywordstyle=\bfseries,
     %
     numbers=left,
     numbersep=6pt,
     numberstyle=\tiny,
     numberfirstline = false,
     breaklines=true,
   }
 
% minimizar fragmentado de listados
\lstnewenvironment{listing}[1][]
   {\lstset{#1}\pagebreak[0]}{\pagebreak[0]}

\lstdefinestyle{CodigoC}
   {
	basicstyle=\scriptsize,
	frame=single,
	language=C,
	numbers=left
   }

\lstdefinestyle{CodigoC++}
   {
	basicstyle=\small,
	frame=single,
	backgroundcolor=\color{gray30},
	language=C++,
	numbers=left
   }
 
\lstdefinestyle{Consola}
   {basicstyle=\scriptsize\bf\ttfamily,
    backgroundcolor=\color{gray30},
    frame=single,
    numbers=none
   }

\newcommand{\bigrule}{\titlerule[0.5mm]}


%Para conseguir que en las páginas en blanco no ponga cabecerass
\makeatletter
\def\clearpage{%
  \ifvmode
    \ifnum \@dbltopnum =\m@ne
      \ifdim \pagetotal <\topskip
        \hbox{}
      \fi
    \fi
  \fi
  \newpage
  \thispagestyle{empty}
  \write\m@ne{}
  \vbox{}
  \penalty -\@Mi
}
\makeatother

\usepackage{pdfpages}
\begin{document}
\begin{titlepage}
 
 
\newlength{\centeroffset}
\setlength{\centeroffset}{-0.5\oddsidemargin}
\addtolength{\centeroffset}{0.5\evensidemargin}
\thispagestyle{empty}

\noindent\hspace*{\centeroffset}

\begin{minipage}{\textwidth}

\centering
\includegraphics[width=0.9\textwidth]{imagenes/logo_ugr.jpg}\\[1.4cm]

\textsc{ \Large \MakeUppercase{\myDocument}\\[0.2cm]}
\textsc{ \MakeUppercase{\myDegree}}\\[1cm]
% Upper part of the page
% 
% Title
{\Huge\bfseries \myTitle \\
}
\noindent\rule[-1ex]{\textwidth}{3pt}\\[3.5ex]
{\large\bfseries \myCompany}\\[0.6cm]
\includegraphics[width=0.5\textwidth]{imagenes/bi4group.png}\\

\end{minipage}

\vspace{1cm}
\noindent\hspace*{\centeroffset}

\begin{minipage}{\textwidth}
\centering

\textbf{Autor}\\ {\myName}\\[2.5ex]
\textbf{Directores}\\
{\myProf\\
\myOtherProf}\\[1cm]
\includegraphics[width=0.3\textwidth]{imagenes/etsiit_logo.png}\\[0.1cm]
\textsc{\myFacultyShort}\\
\textsc{---}\\
\myLocation, \myTime
\end{minipage}
%\addtolength{\textwidth}{\centeroffset}
%\vspace{\stretch{2}}
\end{titlepage}



\frontmatter
\tableofcontents


\chapter{Introducción}

Lala
\chapter{Descripción de la entidad donde realizó las prácticas}

Lala
\chapter{Trabajo realizado}

Lala
\chapter{Valoración personal del trabajo realizado}

El desarrollo de estas prácticas me ha supuesto una experiencia muy interesante y mi primer acercamiento real al mundo empresarial, ya que nunca había trabajado o hecho prácticas como ingeniero informático.

Me ha servido para ver como se trabaja fuera de la universidad, como por ejemplo comprobar desde dentro y con un punto de vista práctico algunas metodologías de desarrollo ágiles que solo había estudiado teóricamente, darme cuenta de como algunas tecnologías aprendidas a lo largo del máster (\textit{Docker}, \textit{ReactJS}, \textit{Flask}, ...) se utilizan realmente o se están implantando en el mundo empresarial, lo que convierte tú formación en un valor adicional, que las empresas valoran.

Centrándome en mi trabajo desarrollado, como es natural al principio estaba un poco perdido, tuve múltiples problemas hasta que comprendí el funcionamiento de las librerías internas de la empresa por la falta de documentación. Una vez empecé a comprender esto mi desarrollo se hizo mucho más fluido, prácticamente al nivel de un miembro del equipo más (llegando incluso a quedarme la última semana solo en el equipo de desarrollo). Por ello puedo decir que estoy contento con mi capacidad de adaptación, ya que en mi escaso período de prácticas he ayudado a desarrollar un proyecto funcional que es actualmente utilizado por la empresa.

Esto mismo ha sido apuntado por algunos de mis responsables directos en la empresa, lo que ha hecho que me hayan tenido en consideración para entrar a formar parte de un nuevo proyecto que comenzará en octubre y me han ofrecido un puesto de trabajo.
\chapter{Conclusiones}

Lala

\bibliography{bibliografia/bibliografia}
\addcontentsline{toc}{chapter}{Bibliografía}
\bibliographystyle{plain}

\end{document}
